\begin{abstract}
Depth-pruned language models are commonly continued-pretrained to restore
aggregate language-model quality, but lower perplexity does not by itself
identify which behaviors have returned. We present a trajectory-level study of
post-pruning recovery on OLMo-2 base models. Prefix variants retain pretrained
lower blocks, append a fresh tail, and continue pretraining; we evaluate
held-out likelihood, two MMLU scoring interfaces, three closed-book QA tasks,
and a broader likelihood suite. In the 16-layer keep14 arm, perplexity improves
to 10.561 after 200k steps but remains $1.428\times$ the intact model's 7.398,
while answer-letter MMLU reaches only .319 versus .605. An intact 32-layer
control stays near the base through its observed 25k plateau checkpoint,
making short-horizon corpus shift an insufficient explanation. Complete-option
scoring raises keep14 to .383, consistent with an answer-interface component;
however, a fully random 16-layer model reaches nearly the same content score
while remaining at chance on answer letters, exposing a large fluency floor.
The deficit also persists on PopQA, TriviaQA, and NQ-open. A different
construction, non-contiguous ShortGPT-16, reaches .474
MMLU at the same 200k target budget, although coupled structural differences
prevent attribution to a single retained block or selection rule. Our conclusion
is methodological: after a structural intervention, likelihood, target
capability, scoring interface, exact construction, and recovery compute should be
reported separately; perplexity measures distributional fit, not a sufficient
certificate of capability recovery.
\end{abstract}
